% Capítulo 4. Revisión bibliográfica y evolución metodológica (rúbrica: revisión 5 pts y parte de estrategia 15 pts).

\section{Revisión crítica de las familias de métodos}
\label{sec:familias}

El problema del capítulo~\ref{ch:problema} pertenece a la familia de programación de vehículos eléctricos (E-VSP). \citeA{perumal_evspreview_2022} ordenan esa literatura según qué decisiones se toman en conjunto, entre ellas asignar corridas a buses, programar la carga, ubicar cargadores y dimensionar la flota. Este caso exige las dos primeras a la vez, con varios depósitos, carga parcial y cargadores finitos, sobre una red de decenas de miles de corridas. Con ese criterio se evaluó cada familia de métodos (Tabla~\ref{tab:familias}).

\begin{table}[htbp]
\centering
\caption{Familias de métodos revisadas, contrastadas con el caso.}
\label{tab:familias}
\footnotesize
\begin{tabularx}{\textwidth}{@{}p{2.3cm} X X p{2.4cm}@{}}
\toprule
Familia & Fortaleza & Debilidad o supuesto que choca con el caso & Referencias \\
\midrule
Flujo en red tiempo-espacio & Exacta y la que mejor escala sin batería; agrega arcos de espera y reduce el modelo. & Con la batería deja de ser un flujo puro; la energía se agrega después o como dimensión del nodo. & \citeA{kliewer_timespacemdvsp_2006}; \citeA{olsen_flowdecomposition_2022} \\
Generación de columnas y \emph{branch-and-price} & Óptimo con cota; trata carga parcial y cargadores finitos. & El mayor caso publicado resuelto al óptimo tiene 101 corridas; la red tiene tres órdenes de magnitud más. & \citeA{devos_capacitatedcharging_2024}; \citeA{wu_waitingtimebranchprice_2025}; \citeA{li_limitedenergy_2014} \\
Metaheurísticas y búsqueda en vecindarios grandes & Escalan a redes reales y entregan planes factibles desde temprano. & No traen cota propia, así que no indican qué tan lejos quedan del óptimo. & \citeA{wen_alnsevsp_2016}; \citeA{limmer_lnsevfleet_2023}; \citeA{ahuja_vlsn_2002} \\
Descomposición de Benders & Separa inversión (flota, cargadores) de operación. & Aquí la flota y la infraestructura vienen dadas y no hay un nivel de inversión que separar. & \citeA{legault_bendersfleetelectrification_2025} \\
Modelos con incertidumbre & Consumo y tiempos aleatorios dentro del modelo. & Los datos no traen distribuciones y el cómputo se multiplica. & \citeA{ricard_chanceconstrained_2026} \\
Simulación & Mide qué tan frágil es un plan. & No optimiza; solo evalúa. & --- \\
\bottomrule
\end{tabularx}
\end{table}

Los métodos exactos garantizan el óptimo, pero no llegan al tamaño de la red. \citeA{wu_waitingtimebranchprice_2025} resuelven al óptimo instancias de hasta 101 corridas y \citeA{devos_capacitatedcharging_2024}, que tratan la variante con cargadores finitos y carga parcial, siguen lejos de las \CorridasLaboral{} corridas del día laboral. Las metaheurísticas sí escalan, y \citeA{wen_alnsevsp_2016} muestran que una búsqueda adaptativa en vecindarios grandes resuelve el E-VSP con varios depósitos y carga parcial. Sin embargo, ninguna de las fuentes revisadas de esa familia modela cargadores finitos ni entrega una cota. Benders no aplica porque no se decide inversión, y para el modelo con incertidumbre faltan distribuciones, así que la robustez, que sigue siendo un tema abierto en el área \cite{wang_bevretrospective_2025}, se estudia con sensibilidad y simulación.

\section{Programar primero y cargar después}
\label{sec:evolucion}

Una estrategia común en la literatura es secuencial. Primero se arman cadenas de corridas como si los buses fueran diésel y después se les agrega la carga \cite{perumal_evspreview_2022}. \citeA{olsen_flowdecomposition_2022} la aplican en redes grandes. Resuelven un flujo en red tiempo-espacio sin batería, lo descomponen en cadenas con reglas que anticipan la energía y recién al final insertan las cargas, agregando buses donde la energía no alcanza. Su ventaja es la escala, y su riesgo es que el flujo no ve la batería, de modo que la fase de carga debe cortar cadenas y agregar buses. \citeA{perumal_evspreview_2022} distinguen los enfoques secuenciales de los integrados, que deciden corridas y carga a la vez y capturan su interacción a cambio de modelos más difíciles, y \citeA{vankootenniekerk_schedulingev_2017} formulan ese problema integrado con carga parcial.

Este proyecto midió ese riesgo en su propio caso. El prototipo de la Entrega~1 siguió la estrategia secuencial de \citeA{olsen_flowdecomposition_2022} sobre la muestra de \EdosRutasMuestra{} rutas. El flujo sin batería, que es una relajación del problema, encadenó las corridas con \EdosFlotaFaseUno{} buses. Cuando se exigió que cada bus tuviera energía para sus corridas, como obliga el problema, la flota subió a \EdosBusesPrototipoSinMejora{} (\EdosAumentoBateria{} buses más, +\EdosPctAumentoBateria\,\%). Ese aumento se explica por haber ignorado la batería al encadenar las corridas. Una construcción que revisa horario y energía en cada asignación logró \EdosBusesConstruccion{} buses factibles con las mismas corridas, de modo que al menos \EdosEvitables{} de los buses agregados se podían evitar reorganizando las jornadas (Tabla~\ref{tab:evolucion}). De ahí sale la decisión central del método, que es considerar la energía al armar cada jornada y no después.

\section{Evolución del método entre las entregas}
\label{sec:evolucion-entregas}

La Tabla~\ref{tab:evolucion} sigue el método a lo largo del semestre. En la Entrega~1 se eligió la red tiempo-espacio por sectores. En la Entrega~2 el diagnóstico anterior llevó a la construcción cronológica con energía y a una mejora local por pares y tríos de jornadas, con un modelo entero que decide las recargas. En la muestra de \EdosRutasMuestra{} rutas logró \EdosBusesModelo{} buses contra \EdosBusesBase{} del caso base (\EdosDeltaPctBuses\,\%). En la Entrega~3 el método se llevó a la red completa y los tres tipos de día, se generalizó la mejora a grupos de 3 a 8 buses elegidos parecidos o al azar, y se agregó un modelo exacto para medir la distancia al óptimo.

\begin{table}[htbp]
\centering
\caption{Evolución del método sobre la misma muestra de \EdosRutasMuestra{} rutas y \EdosCorridasMuestra{} corridas laborales (Entregas 1 y 2) y en la red completa (Entrega 3).}
\label{tab:evolucion}
\small
\begin{tabularx}{\textwidth}{@{}l X r@{}}
\toprule
Entrega & Método & Buses \\
\midrule
1 & Flujo en red tiempo-espacio sin batería (relajación, no operable) & \EdosFlotaFaseUno \\
1 & + descomposición en cadenas y carga insertada después & \EdosBusesPrototipoSinMejora \\
1--2 & + reparaciones y mejora por pares y tríos (prototipo final) & \EdosBusesPrototipo \\
2 & Construcción cronológica con energía, sin mejora & \EdosBusesConstruccion \\
2 & + mejora por pares y tríos con MILP de recargas & \EdosBusesModelo \\
\midrule
3 & Red completa, plan inicial (construcción cronológica) & \BusesInicialLaboral \\
3 & + pares y tríos, 40 min & \BusesMetodoParesTrios \\
3 & + cuartetos al azar, 40 min (método final) & \BusesPlanLaboral \\
\bottomrule
\end{tabularx}
\fuente{Cifras de las Entregas 1 y 2 en \texttt{08-entrega2/informe2/cifras.tex} y de la Entrega 3 en \texttt{09-entrega3/octetos/resultados/red\_completa\_417\_L.csv}.}
\end{table}

\section{Justificación del método elegido}
\label{sec:justificacion}

El método final es una búsqueda en vecindarios grandes con reparación exacta que combina tres ideas con respaldo en la literatura.

La primera es construir con la energía a la vista. El plan inicial recorre las corridas en orden de salida y asigna cada una a un bus que llega a tiempo y tiene energía, como la heurística constructiva con que parte la búsqueda de \citeA{wen_alnsevsp_2016}.

La segunda es mejorar el plan con vecindarios grandes. Una búsqueda local mejora un plan cambiando una parte de él, y \citeA{ahuja_vlsn_2002} describen el dilema de fondo. Un vecindario más grande encuentra mejores cambios en cada intento, pero cada intento cuesta más, y solo conviene si se puede revisar rápido. En el E-VSP, \citeA{wen_alnsevsp_2016} destruyen y reconstruyen partes del plan, y \citeA{limmer_lnsevfleet_2023} reparan la parte destruida con un modelo entero, que es exactamente lo que hace aquí el MILP de recargas.

La tercera es medir el método contra el exacto. Como la búsqueda no trae cota, su calidad se mide resolviendo al óptimo casos chicos con el mismo costo y las mismas reglas, como se propuso desde el informe~1.
